% This file was converted to LaTeX by Writer2LaTeX ver. 1.9.9
% see http://writer2latex.sourceforge.net for more info
\documentclass{article}
\usepackage{calc,amsmath,amssymb,amsfonts}
\usepackage[LGR,T1]{fontenc}
\usepackage[greek,italian]{babel}
\usepackage[style=numeric,backend=biber]{biblatex}
\author{utente}
\date{2026-08-18}
\begin{document}
Prologo

$\text{\textgreek{«}}$Va bene, proviamo.$\text{\textgreek{»}}$

Era il 1983 quando venni a sapere che nella scuola media del mio paese sarebbe stato organizzato un corso pomeridiano di
informatica. Non ricordo chi me ne parlò per primo, né con quali parole mi venne presentato; ricordo, invece, la
sensazione immediata che quella notizia provocò in me. Da qualche tempo sentivo parlare sempre più spesso di computer,
di programmi, di macchine che non si limitavano a eseguire giochi già pronti ma potevano essere istruite, quasi fosse
possibile stabilire con loro una forma di dialogo. Dopo il primo incontro con Pong, la domanda che mi accompagnava da
anni continuava a tornare ogni volta che vedevo un'immagine muoversi su uno schermo: come faceva una macchina a sapere
che cosa doveva fare? Quel corso, pur senza promettermi alcuna risposta precisa, mi sembrò improvvisamente il luogo in
cui avrei potuto cominciare a scoprirlo.

C'era però un ostacolo che, alla mia età, appariva tutt'altro che secondario: frequentavo ancora la quinta elementare,
mentre il corso era destinato agli studenti delle scuole medie. Per qualche giorno cercai di rassegnarmi all'idea che
non avrei potuto partecipare, ma ogni volta che tornavo a pensarci la rinuncia mi sembrava più difficile. Non avevo mai
visto da vicino un vero computer programmabile e non sapevo nemmeno con certezza quale macchina sarebbe stata
utilizzata; proprio per questo l'immaginazione riempiva tutto ciò che ancora non conoscevo, trasformando quell'aula,
prima ancora di entrarci, in una specie di territorio proibito al quale desideravo accedere.

Alla fine ne parlai con i miei genitori. Di quel pranzo non ricordo un discorso preparato né particolari insistenze. Mia
madre ascoltò con attenzione, mentre mio padre, uomo di poche parole, non mostrò alcuna contrarietà. In casa avevano
ormai capito che il mio interesse per la tecnologia non era uno dei tanti entusiasmi destinati a durare il tempo di una
moda. Quando qualcosa mi incuriosiva, non mi bastava usarla: volevo sapere che cosa ci fosse dentro, quale meccanismo
la facesse funzionare e perché, premendo un tasto o spostando una leva, accadesse proprio quella cosa e non un'altra.
Il problema, quindi, non era convincere loro. Bisognava convincere la scuola ad accettare in un corso per ragazzi più
grandi un bambino che non aveva ancora terminato le elementari.

Fu mia madre a chiedere se fosse possibile fare un'eccezione. Non ricordo riunioni, discussioni o attese particolarmente
solenni; il ricordo che mi è rimasto è molto più semplice e proprio per questo, forse, più importante. La risposta
arrivò quasi senza enfasi: $\text{\textgreek{«}}$Va bene, proviamo.$\text{\textgreek{»}}$

Allora mi sembrarono soltanto tre parole capaci di aprirmi la porta di un'aula. Oggi, dopo una vita trascorsa davanti ai
computer, so che aprirono molto di più.

Accanto alla cattedra c'era lui

Il giorno della prima lezione entrai nella scuola media con quella particolare miscela di curiosità e soggezione che si
prova quando ci si trova in un luogo nel quale, per età o per abitudine, si ha la sensazione di non appartenere ancora
del tutto. I corridoi, le porte delle aule, i banchi più grandi di quelli a cui ero abituato contribuivano a ricordarmi
che ero lì grazie a una piccola eccezione, ma bastò arrivare davanti alla stanza del corso perché ogni altra
considerazione perdesse importanza.

Era una normale aula scolastica dei primi anni Ottanta. La lavagna nera occupava la parete di fronte ai banchi, la
cattedra mostrava i segni lasciati da anni di lezioni e nulla, a prima vista, faceva pensare che lì dentro stesse per
accadere qualcosa di diverso da una qualunque attività pomeridiana. Poi lo vidi.

Accanto alla cattedra c'era il Commodore 64.

Per la prima volta quel nome, e soprattutto quel numero che avevo sentito pronunciare tante volte, acquistavano una
forma reale. Osservavo la tastiera beige, i tasti scuri, il logo nell'angolo superiore, il monitor sistemato accanto e,
poco più in là, il drive per i floppy disk, che ai miei occhi possedeva quasi lo stesso fascino del computer. Fino a
quel momento quelle macchine erano appartenute soprattutto ai racconti degli altri, alle frasi che cominciavano con
$\text{\textgreek{«}}$pare che$\text{\textgreek{»}}$ e alle descrizioni inevitabilmente ingigantite dal passaggio da un
ragazzo all'altro. Adesso una di quelle macchine era lì, a pochi metri da me, e potevo finalmente guardarla senza dover
immaginare nulla.

Il Commodore non apparteneva alla scuola. Il professor Rossi Brunori lo portava personalmente ogni volta che teneva
lezione e ad aiutarlo c'era suo figlio, più o meno della nostra età. Li osservai sistemare sulla cattedra il computer,
il monitor e il drive, collegare i cavi con gesti tranquilli e controllare che tutto fosse al proprio posto. Non c'era
un computer per ogni studente, non esisteva una rete, non c'era un laboratorio attrezzato come quelli che molti anni
dopo sarebbero diventati normali perfino nelle scuole. C'era una sola macchina, intorno alla quale ci saremmo raccolti
tutti, eppure nessuno di noi ebbe la sensazione che mancasse qualcosa. Al contrario, bastarono quei pochi oggetti
appoggiati sulla cattedra perché una comune aula scolastica diventasse, almeno ai miei occhi, il luogo più moderno che
avessi mai visto.

Ricordo ancora il modo in cui Rossi Brunori si muoveva accanto al computer. Non aveva l'atteggiamento di chi desidera
dimostrare quanto ne sappia più degli altri e non trasformava la propria competenza in una distanza da colmare. Parlava
con calma, sorrideva spesso e sembrava divertirsi sinceramente nel vedere le nostre reazioni. Forse proprio per questo
gli argomenti nuovi non apparivano mai come ostacoli dietro i quali nascondere la difficoltà della materia: diventavano
domande alle quali provare a dare una risposta insieme.

Molti anni dopo, quando mi è capitato a mia volta di spiegare informatica ad altre persone, ho capito che il suo talento
non consisteva soltanto nel conoscere il BASIC o il Commodore 64. Consisteva nel riuscire a togliere alla tecnologia
quell'aria di mistero che spesso spaventa chi la incontra per la prima volta. Rossi Brunori non ci dimostrava quanto
fosse complicato un computer; ci faceva venire voglia di capire il passaggio successivo. In questo senso, più che
insegnarci il BASIC, ci insegnava a non averne paura.

READY.

Quando tutti ebbero preso posto, il professore si avvicinò alla cattedra e accese il Commodore. Non preparò il momento,
non chiese silenzio e non fece nulla per renderlo solenne. Premette semplicemente l'interruttore. Dopo un brevissimo
istante il monitor si illuminò e comparve quella schermata azzurra, incorniciata dal bordo blu, che negli anni
successivi avrei visto migliaia di volte ma che allora osservavo per la prima volta.

Nelle righe iniziali comparivano parole e numeri che per me non avevano ancora alcun significato preciso: COMMODORE 64
BASIC V2, poi 38911 BASIC BYTES FREE e, infine, una parola brevissima accompagnata dal cursore lampeggiante: READY.

Non ricordo se qualcuno parlò subito. Io certamente no. Guardavo quella schermata cercando di capire che cosa ci fosse
di tanto speciale, e nello stesso tempo avevo la sensazione che ogni elemento mostrato dal computer dovesse essere lì
per un motivo. Rossi Brunori lasciò trascorrere qualche secondo prima di rivolgersi alla classe.

$\text{\textgreek{«}}$Allora, chi di voi mi sa dire perché questo computer si chiama Commodore
Sessantaquattro?$\text{\textgreek{»}}$

Nessuno rispose. Qualcuno abbassò lo sguardo, altri accennarono un sorriso, mentre io provavo inutilmente a immaginare
che cosa potesse significare quel numero. Il professore osservò la classe e, invece di mostrarsi deluso, sembrò quasi
soddisfatto.

$\text{\textgreek{«}}$Perfetto. È proprio la risposta che mi aspettavo.$\text{\textgreek{»}}$

Fu una delle prime cose che imparai del suo modo di insegnare: non sembrava porre domande per scoprire chi sapesse già
la risposta, ma per farci accorgere di non averne ancora una. Indicò il nome del computer sul monitor e ci spiegò che
quel sessantaquattro non era stato scelto soltanto perché suonasse bene; parlava della memoria della macchina. Prima,
però, avremmo dovuto capire che cosa significassero parole come byte e kilobyte, termini che allora ascoltai con la
stessa curiosità con cui avrei ascoltato i nomi degli strumenti di un mestiere completamente sconosciuto.

Il professore non si fermò a una definizione. Indicò il numero 38911 e ci fece notare una stranezza: se il Commodore
possedeva sessantaquattro kilobyte di memoria, perché ne dichiarava disponibili soltanto trentottomilanovecentoundici
byte?

$\text{\textgreek{«}}$Gli altri dove sono finiti? Qualcuno li ha rubati durante la notte?$\text{\textgreek{»}}$

L'aula rise, ma la domanda rimase. Rossi Brunori cominciò a spiegare che non tutta la memoria della macchina poteva
essere utilizzata liberamente dai nostri programmi, perché una parte era già impegnata dal sistema e dall'interprete
BASIC. Non ricordo con certezza tutti gli esempi che utilizzò; ricordo, invece, la sensazione di vedere quel numero
trasformarsi da dettaglio incomprensibile a indizio. Per la prima volta capivo che un computer non mostrava
informazioni a caso: se qualcosa compariva sullo schermo, valeva la pena chiedersi perché.

La lavagna rimase quasi completamente vuota. Rossi Brunori preferiva scrivere direttamente sul Commodore, perché ogni
comando poteva essere visto nello stesso momento in cui veniva spiegato e ogni risultato appariva immediatamente
davanti a noi. Anche gli errori diventavano parte della lezione. Durante una dimostrazione digitò qualcosa in modo
errato e sul monitor comparve il messaggio ?SYNTAX ERROR. Lui lo guardò per un istante e, con la stessa naturalezza con
cui avrebbe commentato una virgola fuori posto, disse: $\text{\textgreek{«}}$Ecco, qui abbiamo fatto un errore di
sbaglio.$\text{\textgreek{»}}$

Qualcuno rise. Anch'io sorrisi, soprattutto per quella strana espressione che sembrava contenere due volte lo stesso
concetto. Rossi Brunori non la spiegò e non cercò l'effetto; corresse ciò che aveva scritto e andò avanti. Proprio per
questo mi rimase impressa. Negli incontri successivi l'avrei sentita tornare ogni volta che qualcosa nel programma non
produceva il risultato previsto, fino a diventare quasi un marchio del suo modo di parlare. Ancora oggi, quando ripenso
a quelle lezioni, è una delle prime frasi con cui riesco a sentirne nuovamente la voce.

Come si parla con un computer?

Più tardi Rossi Brunori tornò a indicare la parola READY. e ci chiese che cosa significasse. Le risposte arrivarono più
facilmente: qualcuno disse che il computer era acceso, qualcun altro che funzionava, un altro ancora che potevamo
iniziare a usarlo. Il professore annuì e ci invitò a prenderla alla lettera: ready significava semplicemente pronto. Il
Commodore aveva terminato ciò che doveva fare all'accensione e adesso aspettava che qualcuno gli impartisse
un'istruzione.

Fino a quel momento il cursore lampeggiante era stato per me poco più di un segno sullo schermo. Dopo quella spiegazione
cominciai a guardarlo in un altro modo. Quella macchina non stava facendo nulla perché non aveva ancora ricevuto nulla
da fare. Era pronta, ma aspettava.

$\text{\textgreek{«}}$E allora$\text{\textgreek{»}}$, domandò Rossi Brunori, $\text{\textgreek{«}}$come si parla con un
computer?$\text{\textgreek{»}}$

$\text{\textgreek{«}}$Con la tastiera$\text{\textgreek{»}}$, rispose qualcuno quasi subito.

Il professore sorrise. $\text{\textgreek{«}}$Io non vi ho chiesto con che cosa. Vi ho chiesto
come.$\text{\textgreek{»}}$

La differenza sembrò sfuggire a tutti e fu proprio da quell'esitazione che costruì la spiegazione. Se fossimo andati in
un Paese del quale non conoscevamo la lingua, disse, avremmo potuto avere una bocca, una voce e tutta la voglia del
mondo di parlare, ma non per questo saremmo riusciti a farci capire. La tastiera era soltanto il mezzo attraverso il
quale inviare qualcosa alla macchina; per comunicare davvero serviva un linguaggio condiviso.

$\text{\textgreek{«}}$Quello che useremo noi si chiama BASIC.$\text{\textgreek{»}}$

Pronunciò quella parola senza alcuna enfasi, come se stesse introducendo uno strumento assolutamente normale. Poi ci
spiegò che il computer non avrebbe cercato di interpretare le nostre intenzioni: se avessimo scritto qualcosa che non
apparteneva al suo linguaggio, si sarebbe limitato a segnalarci l'errore. Non avrebbe immaginato ciò che volevamo dire
e non avrebbe cercato di venirci incontro. A distanza di anni, quella rigidità mi appare come una delle prime lezioni
sulla precisione che abbia mai ricevuto; allora, invece, mi colpì soprattutto l'idea che una macchina potesse capire
qualcosa scritto da noi.

10 PRINT {\textquotedbl}CIAO{\textquotedbl}

Rossi Brunori si sedette davanti alla tastiera e disse che era arrivato il momento di provare. Sullo schermo comparve
lentamente la prima riga di programma che ricordo di aver visto nascere davanti ai miei occhi:

10 PRINT {\textquotedbl}CIAO{\textquotedbl}

Non premette subito RETURN. Ci lasciò osservare quella riga e ci chiese come l'avremmo letta. Per noi era semplicemente
$\text{\textgreek{«}}$dieci, print, ciao$\text{\textgreek{»}}$, ma lui ci invitò a scomporla come avrebbe fatto il
computer: linea dieci, comando PRINT, testo CIAO. Fu una distinzione piccola soltanto in apparenza, perché per la prima
volta quelle parole smettevano di sembrare una formula misteriosa e cominciavano ad avere ruoli differenti all'interno
di una struttura.

Qualcuno gli chiese perché avesse usato proprio il numero dieci. Rossi Brunori spiegò che avremmo potuto iniziare anche
da uno, ma numerare le righe a intervalli lasciava spazio per aggiungerne altre in seguito. Non credo che allora avessi
già compreso il valore di quella precauzione; ricordo però che mi colpì l'idea che un programma dovesse essere pensato
non soltanto per funzionare in quel momento, ma anche per poter essere modificato. Era una forma di ordine che non
avevo mai associato a un gioco o a un computer.

Poi premette RETURN.

La riga scomparve dalla posizione in cui l'avevamo digitata e il cursore tornò a lampeggiare. Sullo schermo non comparve
alcun CIAO. Per qualche secondo ci fu silenzio, finché un ragazzo chiese ciò che probabilmente stavamo pensando tutti:
$\text{\textgreek{«}}$Professore, ma non funziona?$\text{\textgreek{»}}$

Rossi Brunori non si affrettò a correggerlo. $\text{\textgreek{«}}$Perché dite che non funziona?$\text{\textgreek{»}}$

$\text{\textgreek{«}}$Perché non ha scritto CIAO.$\text{\textgreek{»}}$

$\text{\textgreek{«}}$È vero. Ma noi gli abbiamo detto di scriverlo adesso?$\text{\textgreek{»}}$

Fu così che ci spiegò la differenza tra inserire una riga di programma ed eseguirla. Il numero all'inizio della linea
aveva fatto capire al BASIC che quell'istruzione doveva essere conservata in memoria, non eseguita immediatamente. Il
Commodore, quindi, non aveva ignorato ciò che avevamo scritto: lo stava semplicemente trattenendo, in attesa di un
altro comando.

Rossi Brunori digitò tre sole lettere:

RUN

Prima di premere RETURN, però, si fermò. Ci chiese che cosa ci aspettassimo di vedere e ascoltò le ipotesi della classe
senza suggerire quale fosse quella corretta. Qualcuno pensava che CIAO sarebbe comparso al centro dello schermo,
qualcun altro sotto il comando. A un certo punto il professore guardò verso di me.

$\text{\textgreek{«}}$Enzo, tu che cosa ne pensi?$\text{\textgreek{»}}$

Non avevo una risposta tecnica e non volevo inventarla. Guardai la riga che avevamo scritto e dissi semplicemente:
$\text{\textgreek{«}}$Penso che farà quello che gli abbiamo scritto.$\text{\textgreek{»}}$

Rossi Brunori sorrise. $\text{\textgreek{«}}$Questa risposta mi piace. Un computer non fa quello che volevate scrivere e
non fa quello che stavate pensando. Fa quello che gli avete scritto.$\text{\textgreek{»}}$

Poi premette RETURN.

Sul monitor comparve una sola parola: CIAO.

Nessun suono, nessuna animazione, nessun effetto che oggi considereremmo degno di attenzione. Erano soltanto cinque
lettere bianche sullo schermo azzurro, eppure ricordo ancora la sensazione che provai nel guardarle. Quel testo non
apparteneva a un gioco preparato da altri e non faceva parte di una dimostrazione registrata in precedenza. Era
comparso perché una persona aveva scritto un'istruzione e la macchina l'aveva eseguita.

Fu in quel momento che qualcosa cambiò nel modo in cui guardavo il computer. Fino ad allora avevo visto macchine capaci
di fare cose straordinarie, ma sempre già decise da qualcun altro. Adesso avevo assistito, per quanto in una forma
elementare, alla nascita di un comportamento. Prima non c'era nulla; poi qualcuno aveva scritto una riga e il Commodore
aveva fatto esattamente ciò che quella riga gli chiedeva.

Ripensandoci oggi, credo che la prima risposta concreta alla domanda nata anni prima davanti a Pong sia arrivata proprio
in quell'aula. Non sapevo ancora come si muovesse una pallina, come si leggesse un joystick o come si costruisse un
videogioco. Non sapevo nemmeno quanto fosse lunga la strada che separava quelle cinque lettere da tutto ciò che avevo
immaginato. Avevo però compreso il principio dal quale sarebbe partito ogni passo successivo: una macchina poteva
trasformare un'idea in qualcosa di visibile, a patto che qualcuno imparasse a spiegarle quell'idea nel modo giusto.

Uscendo da quella prima lezione non avevo ancora imparato molto BASIC. Avevo imparato qualcosa di più importante: per la
prima volta il computer mi era sembrato meno simile a una scatola magica e più simile a una porta.


\bigskip
\end{document}
